\section{政府信息公开}
\begin{enumerate}
  \item 每年都出题不多，并不复杂，题干只是引子。
  \item 政府信息。\zd{行政机关在履行职责过程中制作 \hz 获取的，以一定形式记录、保存的信息}。
        \begin{enumerate}
          \item 行政机关\zd{行政管理过程}当中。区别于\zd{司法信息、党政信息}。
          \item \zd{制作}。各种法律文件，政府许可、监督、检查形成的笔录。
          \item \zd{获取（保存）}。企业备案的章程。只要在政府那搁着的，放在政府那的。
        \end{enumerate}
  \item 公开。
        \begin{enumerate}
          \item 主体。
                \begin{enumerate}
                  \item \zd{谁制作谁公开，谁保存谁公开}。例外。保存 + 制作 \yjt 一手来源，初次接触。
                  \item 派出机关、派出机构也有公开义务。
                  \item 共同制作的信息。\zd{牵头机关}（找老大）。内部程序上征求其他机关意见，15 日未表态视为同意。
                  \item 被授权组织。
                \end{enumerate}
          \item 范围。
                \begin{enumerate}
                  \item \zd{国家秘密，社会稳定一律不}；商业秘密、个人隐私的原则不。例外。
                        \begin{enumerate}
                          \item 书面征求权利人的意见。
                          \item 同意 \yjt 公开；不同意 \yjt 不公开。\zd{涉及公共利益强行公开}；不表态 \yjt 具体情况具体分析。
                        \end{enumerate}
                  \item 内部信息不公开。\zd{家丑}不可外扬，\al{人事任免}。
                  \item 过程性信息不公开。\al{考试延期延到啥时候}。
                  \item 内部信息 vs 过程性信息。自始至终与你没关/与你有关的。\zd{出题会张冠李戴、偷换概念}。
                \end{enumerate}
        \end{enumerate}
  \item 审查机制。
        \begin{enumerate}
          \item 有保密审查需要报到主管部门/保密机关二次确认。
          \item 程序。
                \begin{enumerate}
                  \item 主动公开（\zd{与大家都有关}）。
                        \begin{enumerate}
                          \item 方式。记者招待会、公众号、报纸、微博、电视等。
                          \item 场所。\zd{档案馆、公共的图书馆、政务服务场所}、信息公告栏、电子显示屏。
                          \item \qd{目录和指南}。
                          \item 文本信息形成 \hz 变更 \zd{20} 日内公开。背。
                        \end{enumerate}
                  \item \zd{依申请公开}。
                        \begin{enumerate}
                          \item 取消了这个实质内容。\scx{要证明它的关联性的要求}。
                          \item \zd{都要提供身份证明}。原则书面，可以口头。期限 20 + 20。
                          \item 收到。\yi 当面提交 \er 邮寄 \yjt 签收 \san 电子方式、平信邮寄 \yjt 双方确认。爱考。
                          \item 答复方式。可以提要求；你要的形式，政府能满足即可。
                          \item 原则不收钱。\zd{申请量大、超过处理范围可收信息处理费（不叫成本费）}。> 200 页，40/页。
                          \item 涉密。涉密部分遮盖，不涉密部份提供。
                          \item 不公开告知 + 说明理由。\zd{可诉}。事实判断。
                          \item 缺材料。一次性告知（\zd{材料补正告知书}），\zd{7} 个工作日内补。
                        \end{enumerate}
                  \item 申请和主动公开的转化。\qd{多人多次}。
                  \item 防止申请人滥用权利的条款。\zd{数量频次超过合理范围；主观目的非法；要求加工、汇总、整理}。
                  \item 信息更证。有权则更正，无权则告知找别的有权主体。
                \end{enumerate}
        \end{enumerate}
  \item 监督和保障。举报、提起诉讼 \hz 复议、主动监察/主动监督、年度报告制度。
\end{enumerate}
